% SHARD-ID: AQM-54-DERIVATIVE-DYNAMICS-F3-EXAMPLES
% SHARD-TITLE: Derivative Dynamics from Kähler Differentials over \(\mathbb F_3\)
% SHARD-SUMMARY: Assembles the derivative, cotangent, tangent, and de Rham definitions needed for finite-ring dynamics.
% SHARD-SUMMARY: Proposes the de Rham CSS test as a calculable stabilizer-selection mechanism, with its extra choices explicit.
% SHARD-SUMMARY: Computes the line-by-line examples \(\mathbb F_3\) and \(\mathbb F_3[\epsilon]/(\epsilon^2)\).
% SHARD-KEYWORDS: derivative, Kahler differentials, cotangent, tangent, de Rham, CSS, stabilizer, F3, dual numbers

\section{Derivative Dynamics from Kähler Differentials over \texorpdfstring{\(\mathbb F_3\)}{F3}}
\label{sec:derivative-dynamics-f3-examples}

\subsection{Status and Source Anchors}

This shard is a checked definition-and-example shard plus one labelled
proposal.  The definitions of derivation, Kähler differentials, their
universal property, the quotient exact sequence, polynomial differentials, and
the de Rham complex are sourced from the Stacks algebra snapshot
\path{references/algebraic_geometry/stacks_project_algebra.tex}, lines
33793--33872, 34103--34120, 34310--34340, and 34388--34524.  Tangent space by
dual numbers and the cotangent-tangent duality are sourced from
\path{references/algebraic_geometry/stacks_project_varieties.tex}, lines
2915--3019.  Residue fields are sourced from
\path{references/algebraic_geometry/stacks_project_algebra.tex}, lines
250--278.  AQM-08 and AQM-09 supply the CSS/stabilizer dictionary, AQM-36
supplies the even Artin-Weyl layer for dual numbers, and AQM-53 supplies the
cotangent parity-shift comparison.  The convention for this shard is
\texttt{CONVENTIONS.md} (ao).

\subsection{Lamport Dependency Skeleton}

\[
\begin{array}{rcl}
D1&:&k=\mathbb F_3,\quad A\text{ is a finite commutative }k\text{-algebra}.\\
D2&:&\operatorname{Der}_k(A,M)\text{ is the }A\text{-module of }
      k\text{-derivations }A\to M.\\
D3&:&\Omega^1_{A/k}\text{ represents }M\mapsto\operatorname{Der}_k(A,M).\\
D4&:&T^*_{A/k,x}=\Omega^1_{A/k}\otimes_A\kappa(x),\quad
      T_{A/k,x}=\operatorname{Hom}_A(\Omega^1_{A/k},\kappa(x)).\\
D5&:&A\xrightarrow{d_0}\Omega^1_{A/k}\xrightarrow{d_1}
      \Omega^2_{A/k}\to\cdots\text{ is the algebraic de Rham complex.}\\
P1&:&\text{The de Rham CSS test puts qudits on }C^1=\Omega^1_{A/k}.\\
L1&:&\Omega^1_{\mathbb F_3/\mathbb F_3}=0.\\
L2&:&\Omega^1_{\mathbb F_3[\epsilon]/(\epsilon^2)/\mathbb F_3}
      \cong \mathbb F_3\,d\epsilon.\\
T1&:&\text{For }A=\mathbb F_3[\epsilon]/(\epsilon^2),\text{ the de Rham CSS}
      \text{ test has one }X\text{-check and no logical qudit.}
\end{array}
\]

\subsection{Definitions D1--D5}

Let \(k\to A\) be a ring map and let \(M\) be an \(A\)-module.  A
\(k\)-derivation is an additive map \(D:A\to M\) satisfying
\[
  D(c)=0\quad(c\in k),\qquad
  D(ab)=aD(b)+bD(a).
\]
The Stacks source denotes the resulting \(A\)-module by
\(\operatorname{Der}_k(A,M)\).

The module of Kähler differentials \(\Omega^1_{A/k}\), with universal
derivation \(d:A\to\Omega^1_{A/k}\), is characterized by the functorial
identity
\[
  \operatorname{Hom}_A(\Omega^1_{A/k},M)
  \cong
  \operatorname{Der}_k(A,M).
\]
Thus \(\Omega^1_{A/k}\) is the algebraic cotangent object before any
quantization choice.

If \(x\in\Spec A\), with residue field \(\kappa(x)\), the cotangent fibre is
\[
  T^*_{A/k,x}=\Omega^1_{A/k}\otimes_A\kappa(x),
\]
and the tangent space is
\[
  T_{A/k,x}=\operatorname{Hom}_A(\Omega^1_{A/k},\kappa(x)).
\]
This is the derivative geometry missing from the bare Weyl-Heisenberg
kinematics: Weyl labels give conjugate translations and phases, while
\(\Omega^1\) records which first-order directions are geometric.

The algebraic de Rham complex is
\[
  A=\Omega^0_{A/k}\xrightarrow{d_0}\Omega^1_{A/k}
  \xrightarrow{d_1}\Omega^2_{A/k}\xrightarrow{d_2}\cdots,
\]
where \(\Omega^i_{A/k}=\wedge_A^i\Omega^1_{A/k}\).  The Stacks construction
gives \(d_{i+1}d_i=0\).  For finite \(k\)-algebras, each term is a finite
\(k\)-vector space in the examples below.

\subsection{Calculable Dynamical Proposals P1}

The first proposal to test is the \emph{de Rham CSS test}.  Choose finite
\(k\)-bases for
\[
  C^0=A,\qquad C^1=\Omega^1_{A/k},\qquad C^2=\Omega^2_{A/k}.
\]
Put qudits on \(C^1\).  With the basis dot pairing, form
\[
  R_X=\operatorname{im}(d_0)\subset C^1,\qquad
  R_Z=\operatorname{im}(d_1^T)\subset C^1.
\]
Since \(d_1d_0=0\), every element of \(R_X\) is orthogonal to every element of
\(R_Z\).  By AQM-08, this is exactly the CSS isotropy condition.  After
choosing phases, this gives commuting stabilizer equations.  This is a
candidate dynamical principle: exact one-forms and coexact one-forms become
local checks, and the remaining logical labels are controlled by the finite
de Rham cohomology
\[
  H^1_{\rm dR}(A/k)=\ker d_1/\operatorname{im}d_0.
\]

This test deliberately does not claim that the ring alone chooses all phases,
a Hamiltonian normalization, or a metric.  Three further calculable branches
are therefore left as proposals:
\[
\begin{array}{ll}
\text{cotangent-Weyl test:} & \text{quantize }T^*_{A/k,x}\text{ or }
  \Omega^1_{A/k}\text{ directly;}\\
\text{Artin-jet comparison:} & \text{compare }\Omega^1\text{ with the even }
  A\text{-Weyl layer of AQM-36;}\\
\text{Hodge/Laplacian test:} & \text{add a pairing or trace density and form }
  d^*d+dd^*.
\end{array}
\]
The last branch is least canonical because \(d^*\) requires extra metric or
pairing data.

\subsection{Example 1: The Field \texorpdfstring{\(A=\mathbb F_3\)}{A=F3} L1}

Let \(A=k=\mathbb F_3\).  A \(k\)-derivation \(D:k\to M\) must annihilate every
element of the image of \(k\).  Since the image is all of \(A\),
\[
  D(0)=D(1)=D(2)=0.
\]
Therefore \(\operatorname{Der}_k(k,M)=0\) for every \(k\)-module \(M\).
By the universal property,
\[
  \operatorname{Hom}_k(\Omega^1_{k/k},M)=0
  \quad\text{for all }M,
\]
so
\[
  \Omega^1_{k/k}=0.
\]
The unique point is \(x=(0)\), with \(\kappa(x)=k\).  Hence
\[
  T^*_{k/k,x}=0\otimes_k k=0,\qquad
  T_{k/k,x}=\operatorname{Hom}_k(0,k)=0.
\]
The de Rham complex begins and ends as
\[
  k\longrightarrow0\longrightarrow0.
\]
The de Rham CSS test has no \(C^1\) site.  Thus a reduced one-point field
contributes no intrinsic derivative stabilizer.  This agrees with AQM-40's
caveat: the ambient Weyl system may exist, but stabilizer data do not come for
free from a geometrically discrete point.

\subsection{Example 2: \texorpdfstring{\(A=\mathbb F_3[\epsilon]/(\epsilon^2)\)}{A=F3[epsilon]/epsilon2} L2--T1}

Let
\[
  A=k[\epsilon]/(\epsilon^2),\qquad k=\mathbb F_3.
\]
Write \(A=k\oplus k\epsilon\), with \(\epsilon^2=0\).  Present \(A\) as
\[
  S/I,\qquad S=k[t],\quad I=(t^2),
\]
and let \(\epsilon\) be the image of \(t\).  The Stacks polynomial-ring lemma
gives
\[
  \Omega^1_{S/k}=S\,dt.
\]
The quotient exact sequence gives
\[
  I/I^2\longrightarrow \Omega^1_{S/k}\otimes_S A
  \longrightarrow \Omega^1_{A/k}\longrightarrow0,
\]
where \(t^2\in I\) maps to
\[
  d(t^2)\otimes1=2t\,dt\otimes1=2\epsilon\,d\epsilon.
\]
In \(\mathbb F_3\), the scalar \(2\) is invertible.  Thus the relation is
\[
  \epsilon\,d\epsilon=0.
\]
Consequently
\[
  \Omega^1_{A/k}\cong A\,d\epsilon/(\epsilon\,d\epsilon)
  \cong (A/(\epsilon))\,d\epsilon
  \cong k\,d\epsilon.
\]
The universal derivation is explicit:
\[
  d(a+b\epsilon)=b\,d\epsilon,\qquad a,b\in k.
\]

The spectrum has one prime, \(\mathfrak m=(\epsilon)\), and its residue field
is
\[
  \kappa(\mathfrak m)=A/\mathfrak m\cong k.
\]
Therefore the cotangent fibre is
\[
  T^*_{A/k,\mathfrak m}
  =
  \Omega^1_{A/k}\otimes_A k
  \cong k\,d\epsilon,
\]
and the tangent space is one-dimensional:
\[
  T_{A/k,\mathfrak m}
  =
  \operatorname{Hom}_A(\Omega^1_{A/k},k)
  \cong k.
\]
The tangent vector \(v\) with \(v(d\epsilon)=1\) corresponds to the derivation
\[
  D_v(a+b\epsilon)=b.
\]

Since \(\Omega^1_{A/k}\) is generated by one element, its exterior square is
zero:
\[
  \Omega^2_{A/k}=0.
\]
Thus the de Rham complex is
\[
  A\xrightarrow{d_0}k\,d\epsilon\longrightarrow0,
  \qquad
  d_0(a+b\epsilon)=b\,d\epsilon.
\]
The kernel is \(k\), the image is all of \(k\,d\epsilon\), and hence
\[
  H^0_{\rm dR}(A/k)=k,\qquad
  H^1_{\rm dR}(A/k)=0.
\]

Now apply the de Rham CSS test.  Choose the basis \(d\epsilon\) of
\(C^1=k\,d\epsilon\).  Then
\[
  R_X=\operatorname{im}d_0=k\,d\epsilon,\qquad
  R_Z=0.
\]
The Pauli label space for the single \(C^1\)-qutrit is \(k\oplus k\), and the
candidate stabilizer label subspace is
\[
  L=(R_X\oplus0)=k\oplus0.
\]
It is isotropic because all \(X\)-type labels commute with each other.  Its
rank is \(1\), so the corresponding all-\(+1\) stabilizer space has dimension
\[
  3^{1-1}=1.
\]
This is the first toy calculation where the nilpotent first-order geometry
selects a nontrivial stabilizer equation.  It is not yet a general dynamical
law; it is a checked finite test case for convention (ao).
